\documentclass[12pt]{article}
\usepackage[margin=1in]{geometry}
\usepackage{amsmath,amssymb,amsfonts}
\usepackage{graphicx}
\usepackage{booktabs}
\usepackage{hyperref}
\usepackage{authblk}
\usepackage{setspace}
\onehalfspacing

\newcommand{\si}[1]{\mathrm{#1}}

\title{Decoupling Safety Guarantees from Policy Competence in\\
Multi-Robot Control Barrier Function Filters}
\author[1]{Harsh Pandhe}
\affil[1]{Independent Research \\ \texttt{harshpandhehome@gmail.com}}
\date{}

\begin{document}
\maketitle

\begin{abstract}
\noindent
Control barrier function (CBF) safety filters are increasingly used to wrap
learned policies in multi-robot systems, on the premise that the filter makes
the resulting system ``safe.'' We argue this premise conflates two distinct
properties: velocity-level safety (the filter never commands a colliding
action) and episode-level competence (the policy accomplishes the task). We
formalize this distinction for a dual-lookahead CBF/ACAS filter applied to a
differential-drive robot, and provide direct empirical evidence of the gap
using a real multi-agent reinforcement learning (MARL) system: a
centralized-training/decentralized-execution (CTDE) MAPPO policy for
cooperative area coverage, trained under a short iteration budget that left
it undertrained. Despite the safety filter being active on every control step
and functioning as designed, the resulting policy still logs a
higher\hspace{0pt}-\hspace{0pt}than\hspace{0pt}-\hspace{0pt}baseline
environment-level collision rate ($3.0$ per episode vs.\ $0.8$--$2.4$ for two
non-learned baselines) and a mean swarm displacement of $0.0\,\si{m}$ --
i.e., the filter successfully keeps a stationary, incompetent policy from
driving into obstacles, but cannot manufacture the exploratory behavior the
policy never learned. We use this as a worked example to propose a simple
three-class taxonomy for characterizing safety-filtered learned systems
(competent-and-filtered, incompetent-and-filtered, unfiltered) and argue that
reported ``collision rate'' figures for CBF-filtered policies should be
accompanied by a competence metric, since the former alone cannot
distinguish a genuinely safe capable system from a safe-by-paralysis one.
\end{abstract}

\noindent\textbf{Keywords:} control barrier functions, safe reinforcement
learning, multi-robot systems, safety filtering, policy evaluation

\section{Introduction}
A control barrier function (CBF) renders a designated safe set forward
invariant by solving, at each control step, a quadratic program (QP) that
minimally perturbs a nominal control input subject to a barrier
constraint~\cite{ames2019cbf}. In multi-robot and learning-enabled robotics,
CBF filters are frequently proposed as a way to certify the safety of an
otherwise uncertified learned policy: the policy proposes an action, the
filter either passes it through or overrides it, and the composed system
inherits a formal collision-avoidance guarantee independent of the policy's
internal quality~\cite{cheng2019end}. Recent work continues to extend this
pattern -- enforcing CBF constraints \emph{during} RL training rather than
only at
deployment~\cite{cbfrl2025}, combining CBFs with actor-critic architectures
via learned Koopman
representations~\cite{koopmancbf2026}, and applying CBF-informed RL to
connected and automated vehicles~\cite{beyondsafety2026} and mobile robot
navigation under
risk~\cite{riskaware2026}.

This literature is, almost uniformly, framed around a positive claim: the
filter makes the system safe. We do not dispute that claim at the level it is
made -- a correctly implemented CBF filter genuinely does bound the
velocity-level action space to avoid detected collisions. Our contribution is
narrower and, we argue, underexamined: we characterize what that guarantee
does \emph{not} imply, using a real, reproducible multi-agent system in which
an undertrained policy remained ``safe'' under an active CBF filter while
being functionally useless. We propose that this distinction -- between
velocity-level safety and episode-level competence -- deserves explicit
separation in how CBF-filtered learned systems are evaluated and reported,
and offer a simple taxonomy for doing so.

\section{Related Work}
Ames et al.~\cite{ames2019cbf} establish the theoretical foundation for CBFs
as a general safety-certification tool via forward invariance of a safe set.
Cheng et al.~\cite{cheng2019end} combine CBFs with RL by using the filter as
an online safety layer around a learned policy for safety-critical continuous
control, similar in spirit to the architecture we analyze; their emphasis,
as in most follow-on work, is on demonstrating that the filter prevents
constraint violation, not on characterizing the space of behaviors the
filtered system can still exhibit when the underlying policy is weak.
CBF-RL~\cite{cbfrl2025} internalizes CBF constraints into the training
process itself so that safety need not be enforced online at deployment; this
is a complementary direction to our analysis, which is specifically about
systems that \emph{do} retain an online filter, since online filtering
remains common practice when training-time guarantees are unavailable or
insufficient. Koopman-based CBF filters for actor-critic
RL~\cite{koopmancbf2026} and CBF-informed RL for connected and automated
vehicles~\cite{beyondsafety2026} extend the filtering paradigm to new system
classes and representations without revisiting the competence question. A
risk-aware dynamic safety filter for mobile robot
navigation~\cite{riskaware2026} is closer to our concern in that it
explicitly reasons about risk rather than binary safety, but does not report
a competence metric alongside its safety metric. Yu et al.~\cite{yu2022mappo}
establish MAPPO with a centralized critic as a strong CTDE baseline for
cooperative multi-agent tasks; we use an architecture in this family as the
learned policy under test, not as a contribution in itself.

\section{System Under Test}
\label{sec:system}
We use a three-robot TurtleBot3 Waffle swarm performing cooperative area
coverage in a Gazebo Sim environment (ROS~2 Jazzy), with a shared MAPPO
policy trained via a centralized critic and a permutation-invariant
neighbor-embedding actor. Full architectural detail is given in a companion
system paper; here we summarize only what is needed to interpret the safety
analysis.

\textbf{Dual-lookahead CBF/ACAS filter.} Given a nominal command
$(v_{\text{nom}}, \omega_{\text{nom}})$ from the policy, the filter solves
\begin{equation}
\min_{v,\omega}\ (v-v_{\text{nom}})^2 + 0.05(\omega-\omega_{\text{nom}})^2
\ \text{s.t.}\ \{A_k v + B_k\omega + \gamma h_k \ge 0\}_k,
\end{equation}
over barrier constraints $h_k$ derived from LiDAR returns and neighbor
positions at two lookahead points ($\pm l$, $l=0.12\,\si{m}$) along the
robot's heading axis, with safety margin $d_{\text{safe}}=0.20\,\si{m}$ and
class-$\mathcal{K}$ gain $\gamma=2.0$, considering all LiDAR sectors and
neighbors within $0.6\,\si{m}$. The QP is solved with SLSQP at every control
step. On infeasibility, the filter halts forward motion while permitting
rotation, rather than passing the unfiltered nominal command through. An
auxiliary Active Collision Avoidance Safeguard (ACAS) commands a slow reverse
whenever raw minimum LiDAR range drops below $0.22\,\si{m}$ with positive
commanded forward speed. This filter is applied identically regardless of
which policy (learned or scripted) produced the nominal command, which is
what makes it possible to isolate the filter's contribution from the
policy's.

\textbf{Environment-level collision metric.} Following standard practice, we
report a binary collision event whenever any robot's minimum LiDAR range
drops below $0.20\,\si{m}$ during an episode -- the same threshold used as
$d_{\text{safe}}$ in the filter. This is deliberately the metric most
favorable to a ``the filter makes it safe'' narrative, since it is measured
at the same distance the filter itself defends.

\section{Empirical Characterization}
\label{sec:empirical}
We benchmark three methods under identical filter conditions in the nominal
regime (5 episodes each): a Random Walk baseline, a reactive Frontier
Heuristic baseline (closest-unvisited-cell target selection with reactive
obstacle avoidance), and the MAPPO policy trained for 15 iterations
($4{,}500$ total environment steps -- physics-in-the-loop training in Gazebo
has no fast-simulator shortcut, so this is a genuinely small budget by MARL
standards). All three methods are subject to the identical CBF/ACAS filter
described in Section~\ref{sec:system}.

\begin{table}[h]
\centering
\caption{Nominal-regime comparison, mean $\pm$ std over 5 episodes. All
methods run under the identical CBF/ACAS filter.}
\label{tab:main}
\begin{tabular}{@{}lcccc@{}}
\toprule
Method & ACR (\%) $\uparrow$ & Redund.\ $\downarrow$ & Dist.\ (m) & Coll.\ $\downarrow$ \\
\midrule
Random Walk        & $18.1\pm0.6$ & 52.64 & 6.3 & 0.80 \\
Frontier Heuristic & $21.1\pm9.9$ & 8.12  & 4.8 & 2.40 \\
MAPPO (undertrained) & $10.4\pm1.5$ & 14.63 & \textbf{0.0} & \textbf{3.00} \\
\bottomrule
\end{tabular}
\end{table}

The result central to this paper's argument is the MAPPO row's
\emph{mean displacement of $0.0\,\si{m}$} co-occurring with its
\emph{highest} collision count of the three methods. Evaluated
deterministically (no exploration noise), the policy collapses to a
near-stationary behavior at or immediately after spawn -- it does not explore
the environment at all -- while still logging more environment-level
collision events than either scripted baseline, which do move. Diagnostic
data from the training run corroborates the underlying cause: policy entropy
fell to $-0.53$ by the final iteration (from a positive value earlier in
training), and stochastic training rollouts had explored substantially more
of the state space than the deterministic evaluation policy retained --
consistent with a policy that converged to a locally ``safe'' near-stationary
behavior before it had seen enough episodes to learn that exploration is
rewarded.

Critically, this is not a filter failure. The filter is not designed to,
and does not claim to, prevent the environment-level collision event that
triggers when a stationary robot's already-close neighbor or wall geometry
places it under the $0.20\,\si{m}$ threshold at spawn, nor to invent an
escape trajectory for a policy that never proposes one. The filter's actual
guarantee -- bounded velocity given current sensing, with a rotation-only
fallback on QP infeasibility rather than passing an unfiltered command -- was
exercised correctly throughout. What the result demonstrates is that this
guarantee, though real, is insufficient on its own to produce the intuitive
notion of a ``safe robot'': a robot that is safe because it is frozen in
place, adjacent to an obstacle it was placed near at spawn, is safe by the
filter's contract and useless by every other measure.

\section{A Taxonomy for Reporting Filtered Systems}
We propose classifying a CBF-filtered learned system along two independent
axes rather than one: \emph{filter status} (active/inactive, and whether it
degrades gracefully on infeasibility) and \emph{policy competence}
(task-level performance under the filter, e.g.\ coverage rate, goal-reach
rate, or an equivalent domain metric). This yields three practically
distinguishable classes:

\begin{itemize}
  \item \textbf{Class A -- Competent and filtered.} The policy accomplishes
        the task and the filter is rarely or never binding. The intuitive
        ``safe and capable'' system.
  \item \textbf{Class B -- Incompetent and filtered.} The policy fails to
        accomplish the task (low task-metric score, as in our
        $0.0\,\si{m}$-displacement MAPPO run) while the filter prevents
        velocity-level constraint violation. Safe-by-paralysis or
        safe-by-restriction: the collision-rate metric alone cannot
        distinguish this from Class A.
  \item \textbf{Class C -- Unfiltered.} No safety filter; task competence and
        safety are entangled, and any observed low collision rate is
        attributable to the policy alone.
\end{itemize}

Under this taxonomy, our MAPPO result is a documented, empirically produced
Class~B instance. We argue that any paper (including a system paper we have
written using the same underlying platform) reporting a collision-rate
figure for a CBF-filtered learned policy should also report a task-competence
figure alongside it, precisely because Table~\ref{tab:main}'s collision
column, read in isolation, does not distinguish MAPPO's Class~B result from
a Class~A one -- the two are visually similar (low collision counts) but
substantively opposite.

\section{Discussion and Limitations}
Our empirical evidence comes from a single system and a single undertrained
policy; we do not claim the entropy-collapse mechanism that produced this
specific Class~B instance is the only route to one. We also do not claim CBF
filters are ineffective -- the filter in Section~\ref{sec:system} performed
exactly as specified throughout every experiment reported here, including the
Class~B case, and a policy-agnostic bounded-velocity guarantee is a genuinely
useful property independent of the concerns raised in this paper. Our claim
is narrower: that the guarantee and the outcome most readers associate with
it (a capable, safe robot) are separable, that this separation is easy to
produce by accident (an undertrained policy, not an adversarial one, was
sufficient here), and that current reporting conventions in the CBF-filtered
RL literature do not consistently surface it. A natural extension is a
systematic sweep across training budgets, reward designs, and filter
parameters to map the boundary between Class~A and Class~B behavior, rather
than the single worked example provided here.

\section{Conclusion}
We characterized, with direct empirical evidence from a real multi-robot
reinforcement learning system, the gap between what a CBF/ACAS safety filter
guarantees (bounded, collision-avoiding velocity given current sensing) and
what it is often implicitly credited with (a competent, safe robot). An
undertrained MAPPO policy remained ``safe'' by the filter's own contract
while being entirely stationary and, by a strict environment-level metric,
more collision-prone than two non-learned baselines that actually moved. We
propose reporting task-competence alongside collision-rate for CBF-filtered
learned systems, and offer a three-class taxonomy (competent-filtered,
incompetent-filtered, unfiltered) as a starting point for making this
distinction explicit in future work.

\bibliographystyle{plain}
\bibliography{topic2_references}

\end{document}
